\section{The Acz\'el Development}
\label{app:aczel}

This appendix sketches the from-scratch construction behind the
product-rule half of Theorem~\ref{thm:cox}.  The body stands without
it.  The target shape is Acz\'el's associativity
theorem~\cite{aczel1966}: a suitably regular associative
$F : \mathbb{R} \to \mathbb{R} \to \mathbb{R}$ regraduates to
multiplication, $G(F\,x\,y) = G(x) \cdot G(y)$ for a strictly
monotone $G$.  \textsf{mathlib} offers no route.  It has no
associativity representation, no H\"older-style embedding of ordered
semigroups on intervals, and the Cauchy functional equation only
under continuity, not monotonicity.

\paragraph{Hypotheses.}
A \texttt{Scale} bundles: associativity of $F$, \emph{divisibility}
$\forall c\, \exists x,\ F\,x\,x = c$, positivity on the cone
$x < F\,x\,c$, joint order-preservation, joint continuity, a unit
$u$, and identity at $-\infty$, meaning
$F\,x\,z \to x$ as $z \to -\infty$.  Order-preservation was absent
from our first bundle, and the additivity proof needs it, a
definition-check finding (Section~\ref{sec:engineering}).
Divisibility is taken as a hypothesis in the bundle, although its
interval form is a consequence of the intermediate value theorem
(\texttt{exists\_diag\_eq}).  Non-vacuity: \emph{log-sum-exp},
$F\,x\,y = \log(e^x + e^y)$ with $u = 0$, satisfies every field
(\texttt{logSumExpScale}), and for it the intended generator is
globally $\exp$.

\paragraph{Construction.}
Write $a^{(n)}$ for $n$ right-associated copies of $a$ under $F$
(\texttt{Fpow}).  Associativity alone makes
$n \mapsto a^{(n)}$ a semigroup morphism from $(\mathbb{N}, +)$, with
no commutativity needed.  Commutativity is a \emph{corollary} of the
generator, not a prerequisite.  One fact usually assumed is here
\emph{derived}: unboundedness of the iterates.  Were they bounded,
monotone convergence would give a fixed point $F\,L\,a = L$,
contradicting positivity, and unboundedness yields the Archimedean
property.  Divisibility supplies a tower of $2^n$-th roots
$u \succ u^{1/2} \succ u^{1/4} \succ \cdots$ (\texttt{droot}) with
the coherence law that $2^\ell$ fine copies recompose one coarse
mark.  Counting how many copies of the $n$-th root fit under $b$
(\texttt{dcount}, well defined by the Archimedean property) gives
dyadic approximants $g_n(b) = \mathrm{dcount}(b,n)/2^n$.  The
refinement sandwich $2N + 1 \le N' \le 2N + 2$ makes them monotone
and Cauchy, and the generator is their limit,
$g(b) = \bigsqcup_n g_n(b)$.

\paragraph{Results (all \texttt{sorry}-free, standard axioms).}
On the cone $\{x \mid u \le x\}$ we prove three things.
Normalization: $g(u) = 1$.  Additivity:
$g(F\,x\,y) = g(x) + g(y)$, by squeezing super- and sub-additive
counting bounds ($\mathrm{dcount}\,x + \mathrm{dcount}\,y + 1 \le
\mathrm{dcount}(F\,x\,y) \le \mathrm{dcount}\,x +
\mathrm{dcount}\,y + 3$) against a vanishing $2^{-n}$ slack.  And
\emph{strict} monotonicity, proved without any continuity of $g$ by
a count-amplification argument: a shrinking root inserted between
$x < y$ contributes a full extra block of $2^{n-M}$ copies.
Exponentiating gives the multiplicative form, a strictly monotone
$G$ with $G(F\,x\,y) = G(x)\,G(y)$ on the cone
(\texttt{exists\_mul\_generator}), which is the conclusion shape of
Acz\'el's theorem restricted to the cone
(\texttt{aczelStatement\_cone}).  A bounded reorientation
$\bar G = \exp(-g)$ transports it to the $[0,1]$, conjunctive
orientation of Section~\ref{sec:cox}.

\paragraph{Honest residue.}
The constructed $g$ is discontinuous at the unit (it is $0$ below
$u$), so matching the classical statement verbatim still needs the
off-cone extension (a group completion) and continuity of the
extension.  That is genuine analysis, open here, but
\emph{witnessed}: for the log-sum-exp archetype the global
continuous generator exists
(\texttt{hasOrderedGenerator\_logSumExp}).  Of H\"older's classical
hypotheses, the Archimedean property is a theorem in this
development rather than an assumption, and the interval form of
divisibility follows from the intermediate value theorem.  Only the
embedding's global regularity remains.
